\clearpage\SourceLeaf{253}
\par\bigskip\Needspace{9\baselineskip}
\pdfbookmark[0]{Weitere Entwicklung der Untersuchungen über reine Gleichungen}{aequationes-purae}
\begin{center}\large\bfseries WEITERE ENTWICKLUNG DER UNTERSUCHUNGEN\par\medskip ÜBER REINE GLEICHUNGEN.\end{center}
\BookRule{.12\linewidth}
\renewcommand{\BookRunningTitle}{Weitere Untersuchungen über reine Gleichungen}
\renewcommand{\Article}[1]{\par\medskip\Needspace{4\baselineskip}\pdfbookmark[1]{Art. #1}{purae-art-#1}\centerline{\textbf{#1.}}\nobreak\smallskip}

\Article{1}
Das Verfahren, mit dem wir in Artikel~360 der \emph{Disquiss.\ Arithm.} gezeigt haben, wie die Gleichung \mbox{$x^n-1=0$} zu lösen ist, bildet eine überaus fruchtbare und bedeutende Theorie, von der wir in jenem Werk nur die ersten Grundlagen berühren konnten. Wir hoffen daher, dass es den Mathematikern willkommen sein wird, wenn wir diesen Gegenstand hier erneut aufnehmen, das nur kurz und zum Teil unter Auslassung der Beweise Skizzierte ausführlicher behandeln und die seit jener Zeit hinzugekommenen Entwicklungen eingehender verfolgen.

Der Exponent $n$ wird als Primzahl vorausgesetzt und die Zahl $n-1$ als in die Faktoren $\alpha\times\beta\times\gamma$ zerlegt; ferner bezeichnen wir mit $g$ irgendeine primitive Wurzel zum Modul $n$. Es bezeichne $r$ eine beliebige Wurzel der Gleichung \mbox{$x^n-1=0$} und $R$ eine beliebige Wurzel der Gleichung \mbox{$x^\beta-1=0$}. Bezeichnen wir also den Umfang des Kreises, dessen Radius $=1$ ist, mit $P$ und die imaginäre Größe $\mathord{\surd}{-1}$ mit $i$, so werden sämtliche Wurzeln der Gleichung \mbox{$x^\beta-1=0$} oder sämtliche Werte von $R$ durch die Formel
\[
\cos\frac{kP}{\theta}+i\sin\frac{kP}{\theta}
\]
dargestellt, wobei $k$ die ganzen Zahlen $0,1,2,3\ldots\beta-1$ durchläuft. Ferner sind offenbar sämtliche Potenzen jeder Wurzel $R$ selbst ebenfalls Wurzeln; ist $R$ die Wurzel, die einem Wert von $k$ entspricht, der zu $\beta$ teilerfremd ist, so sind alle Potenzen $R^0,R,R^2,R^3\ldots R^{\beta-1}$ voneinander verschieden und erschöpfen daher die gesamte Menge der Wurzeln. In diesem Fall nennen wir $R$ eine \emph{eigentliche} Wurzel der Gleichung \mbox{$x^\beta-1=0$}; dagegen heißt eine Wurzel $R$,
\SourceLeaf{254}
die einem Wert von $k$ entspricht, der zu $\beta$ nicht teilerfremd ist, \emph{uneigentlich}. Ist $\delta$ der größte gemeinsame Teiler der Zahlen $k$ und $\beta$, so gilt, wie man leicht sieht, $R^{\frac\beta\delta .}=1$, während alle Potenzen $R^0,R,R^2,R^3\ldots R^{\frac\beta\delta-1}$ voneinander verschieden sind; daher ist $R$ eine eigentliche Wurzel der Gleichung \mbox{$x^{\frac\beta\delta}-1=0$}. Dasselbe gilt für die Gleichung \mbox{$x^n-1=0$}, doch sind deren Wurzeln mit Ausnahme der Wurzel $1$ notwendigerweise sämtlich eigentlich.

\Needspace{20\baselineskip}\Article{2}
Nach diesen Vorbemerkungen soll sich unsere Untersuchung vor allem mit Funktionen der folgenden Form befassen, die aus $\beta\gamma$ Gliedern bestehen:
\[
r+Rr^{g^\alpha}+R^2r^{g^{2\alpha}}+R^3r^{g^{3\alpha}}\ldots+R^{\beta\gamma-1}r^{g^{\alpha\beta\gamma-\alpha}}
\]
Der Kürze halber bezeichnen wir sie mit dem Zeichen $[r,R]$. Die einzelnen Glieder eines solchen Ausdrucks sind Produkte aus Potenzen von $r$ und Potenzen von $R$; die Exponenten der ersteren bilden eine geometrische, die der letzteren eine arithmetische Folge. Die Exponenten
\[
1,\quad g^\alpha,\quad g^{2\alpha},\quad g^{3\alpha}\ldots\ldots g^{\alpha\beta\gamma-\alpha}
\]
sind sämtlich zum Modul $n$ inkongruent zueinander, also jene Potenzen von $r$ voneinander verschieden; weiter fortgesetzt würden sie jedoch dieselbe Reihe erneut beginnen, da $g^{\alpha\beta\gamma}\equiv1(\mathrm{mod.}\,n)$ und daher $r^{g^{\alpha\beta\gamma}}=r$ ist. Die anderen Faktoren
\[
1,\quad R,\quad R^2,\quad R^3\dotsfour R^{\beta\gamma-1}
\]
bilden $\gamma$ gleiche Perioden, da $R^\beta=1$, $R^{\beta-1}=R$ usw. ist. Daraus folgt, dass sich die Funktion $[r,R]$ auch so darstellen lässt:
\[
\begin{aligned}
&r+r^{g^{\alpha\beta}}+r^{g^{2\alpha\beta}}\ldots\ldots+r^{g^{\alpha\beta\gamma-\alpha\beta}}\\
&{}+R\bigl(r^{g^\alpha}+r^{g^{\alpha\beta+\alpha}}+r^{g^{2\alpha\beta+\alpha}}\ldots\ldots+r^{g^{\alpha\beta\gamma-\alpha\beta+\alpha}}\bigr)\\
&{}+R^2\bigl(r^{g^{2\alpha}}+r^{g^{\alpha\beta+2\alpha}}+r^{g^{2\alpha\beta+2\alpha}}\ldots\ldots+r^{g^{\alpha\beta\gamma-\alpha\beta+2\alpha}}\bigr)\\
&{}+\text{usw.}\\
&{}+R^{\beta-1}\bigl(r^{g^{\alpha\beta-\alpha}}+r^{g^{2\alpha\beta-\alpha}}+r^{g^{3\alpha\beta-\alpha}}\ldots\ldots+r^{g^{\alpha\beta\gamma-\alpha}}\bigr).
\end{aligned}
\]
oder, wenn wir das Zeichen aus Artikel~343 der Disq.\ Ar.\ einführen,
\[
[r,R]=(\gamma,1)+R(\gamma,g^\alpha)+R^2(\gamma,g^{2\alpha})\ldots+R^{\beta-1}(\gamma,g^{\alpha\beta-\alpha})
\]

\SourceLeaf{255}
\Needspace{20\baselineskip}\Article{3}
Nehmen wir für die Wurzel $r$ die Einheit, so erhalten wir
\[
[1,R]=1+R+R^2+R^3\ldots+R^{\beta\gamma-1}=\gamma(1+R+R^2+R^3\ldots+R^{\beta-1})
\]
mit dem Wert $=\beta\gamma$, wenn wir auch für $R$ die Wurzel $1$ nehmen, dagegen $=0$ für jeden anderen Wert von $R$. Bleibt hingegen $r$ unbestimmt und setzen wir $R=1$, so ist $[r,1]=r+r^{g^\alpha}+r^{g^{2\alpha}}+r^{g^{3\alpha}}\ldots+r^{g^{\alpha\beta\gamma-\alpha}}$ oder, mit dem in den Disq.\ Ar.\ eingeführten Zeichen, $[r,1]=(\beta\gamma,1)$; das heißt, der Ausdruck besteht aus einer Periode von $\beta\gamma$ Wurzeln, von denen eine $r$ selbst ist. Ist $\alpha=1$, so umfasst diese Periode sämtliche Wurzeln $r,r^2,r^3\ldots r^{n-1}$, nur in anderer Reihenfolge.

Man beachte noch die folgenden Beziehungen, deren Grund unmittelbar ersichtlich ist:
\[
[r,R]=R[r^{g^\alpha},R]=R^2[r^{g^{2\alpha}},R]\quad\text{oder allgemein}\quad=R^k[r^{g^{\alpha k}},R]
\]
wobei $k$ eine beliebige positive ganze Zahl bezeichnet. Daraus folgt, dass die Funktion $[r^m,R]$ entweder $=[1,R]$ ist, nämlich wenn $m$ durch $n$ teilbar ist, oder sich in den übrigen Fällen auf die Form $R^\mu[r^{g^\nu},R]$ mit $\nu<\alpha$ bringen lässt. Ist nämlich $m$ nicht durch $n$ teilbar, so ist es zum Modul $n$ irgendeiner Potenz von $g$ kongruent, deren Exponent sich wie in den Disq.\ Ar.\ zweckmäßig durch $\mathrm{ind.}\,m$ ausdrücken lässt; setzen wir also $\mathrm{ind.}\,m=\lambda\alpha+\nu$, was offenbar so geschehen kann, dass $\nu<\alpha$ ist, so erhalten wir $[r^m,R]=[r^{g^{\lambda\alpha+\nu}},R]=R^{-\lambda}[r^{g^\nu},R]$. Daher ist $\mu=-\lambda$ zu setzen oder, wenn du einen positiven Exponenten wünschst, $\mu\equiv-\lambda(\mathrm{mod.}\,\beta)$.

\Article{4}
\textsc{Satz.} \emph{Bezeichnet $r'$ ebenso wie $r$ eine beliebige Wurzel der Gleichung \mbox{$x^n-1=0$} und $R'$ ebenso wie $R$ eine beliebige Wurzel der Gleichung \mbox{$x^\beta-1=0$}, so ist das Produkt}
\[
\begin{aligned}
&[r,R]\times[r',R']=\\
&\quad[rr',RR']+R[r^{g^\alpha}r',RR']+R^2[r^{g^{2\alpha}}r',RR']\\
&\quad{}+R^3[r^{g^{3\alpha}}r',RR']\ldots+R^{\beta\gamma-1}[r^{g^{\alpha\beta\gamma-\alpha}}r',RR']
\end{aligned}
\]

\emph{Beweis.} Multipliziert man $[r,R]$ mit den einzelnen Gliedern von $[r',R']$, so lässt sich das Produkt in folgender Form darstellen:
\[
\begin{aligned}
&[r,R]r'+RR'[r^{g^\alpha},R]r'^{g^\alpha}+R^2R'^2[r^{g^{2\alpha}},R]r'^{g^{2\alpha}}\\
&{}+R^3R'^3[r^{g^{3\alpha}},R]r'^{g^{3\alpha}}\ldots\\
&{}+R^{\beta\gamma-1}R'^{\beta\gamma-1}[r^{g^{\alpha\beta\gamma-\alpha}},R]r'^{g^{\alpha\beta\gamma-\alpha}}
\end{aligned}
\]

\SourceLeaf{256}
Fasst man dann die ersten Glieder der einzelnen, ordnungsgemäß entwickelten Teile zusammen, so ergibt sich $[rr',RR']$; sammelt man ebenso die zweiten Glieder, so entsteht $R[r^{g^\alpha}r',RR']$, und so weiter. Auf diese Weise erhält man schließlich die angegebene Form des Produkts. Q.\ E.\ D.

Ferner zeigt schon die Vertauschung von $r,R$ mit $r',R'$, dass sich dasselbe Produkt auch in folgende Form bringen lässt:
\[
\begin{aligned}
&[rr',RR']+R'[rr'^{g^\alpha},RR']+R'^2[rr'^{g^{2\alpha}},RR']\\
&{}+R'^3[rr'^{g^{3\alpha}},RR']\dotsfour+R'^{\beta\gamma-1}[rr'^{g^{\alpha\beta\gamma-\alpha}},RR']
\end{aligned}
\]

\Needspace{16\baselineskip}Daraus folgt weiter: Bezeichnen auch $r'',r'''$ usw.\ beliebige Wurzeln der Gleichung \mbox{$x^n-1=0$} sowie $R'',R'''$ usw.\ beliebige Wurzeln der Gleichung \mbox{$x^\beta-1=0$}, so ist das Produkt der Funktionen $[r,R],[r',R'],[r'',R''],[r''',R''']$ usw., wie groß auch ihre Anzahl sein mag, gleich der Summe
\[
\begin{gathered}
\Sigma R'^{k'}R''^{k''}R'''^{k'''}\text{ usw.}\\
{}[rr'^{g^{\alpha k'}}r''^{g^{\alpha k''}}r'''^{g^{\alpha k'''}}\text{ usw.},\ RR'R''R'''\text{ usw.}]
\end{gathered}
\]
wobei für $k',k'',k'''$ usw.\ sämtliche Zahlen $0,1,2,3\ldots\beta\gamma-1$ eingesetzt werden, die auf alle möglichen verschiedenen Arten miteinander kombiniert werden. Dadurch ergeben sich insgesamt $\beta^{\mu-1}\gamma^{\mu-1}$ Glieder, wenn $\mu$ die Anzahl jener miteinander multiplizierten Funktionen bezeichnet.
